\chapter{Using the Scanner}\label{ch:scanner}

The scanner lets your program send requests to a PLC or a device. All
scanner calls block until they are complete or until their timeout. On an
RTOS, call them from a task of their own (section~\ref{sec:threads}).

\section{Sessions}\label{sec:sessions}

A \term{session} is a TCP connection to port 44818 with a registered session
handle. All explicit messages and the Forward\_Open of a Class~1 connection
go through a session.

\subsection{The Slot}

The \code{slot} field of \code{EipOpenConfig} tells the scanner how to send
each request:

\begin{itemize}
  \item \code{slot = -1} sends the request directly to the device. Use it for
        CompactLogix, Micro800, and all devices that you reach directly.
  \item \code{slot} 0--255 puts the request in an Unconnected Send. The
        request goes across ControlLogix backplane port 1 to that slot.
\end{itemize}

\begin{caution}
A zeroed \code{EipOpenConfig} has slot 0, which routes across a backplane.
For a device that you reach directly, set \code{slot = -1}. If you do not,
the device can refuse the requests, or a router can send them to the wrong
module.
\end{caution}

\subsection{Open and Close a Session}

\begin{procedure}{Open a session and read a tag}
\begin{enumerate}
  \item Declare the \code{EipSession} as \code{static}. It is approximately
        4.7~KB.
  \item Set an \code{EipOpenConfig} to zero.
  \item Set \code{ip} to the address of the device, as dotted text.
  \item Set \code{slot}. Use $-1$ for a device that you reach directly.
  \item Optional: set \code{timeout_ms}. The value 0 selects 5000~ms.
  \item Call \code{eip_session_open}.
  \item If the result is not \code{EIP_OK}, stop. Use
        \code{eip_status_text} to show the cause.
  \item Send your requests.
  \item Call \code{eip_session_close}.
\end{enumerate}
\end{procedure}

\begin{lstlisting}
#include <stdio.h>
#include "eip.h"

static EipSession s;   /* about 4.7 KB: keep it off the stack */

int read_speed(void)
{
    EipOpenConfig open = {0};
    uint8_t buf[4];
    uint16_t len = sizeof buf, type = 0;
    uint8_t status = 0;
    int rc;

    open.ip = "192.168.1.10";
    open.slot = -1;           /* direct. A zeroed config means slot 0! */
    open.timeout_ms = 2000;   /* 0 selects 5000 */
    open.vendor_id = 65000;   /* used only by a Class 1 Forward_Open */

    rc = eip_session_open(&s, &open);
    if (rc != EIP_OK) {
        printf("open failed: %s\n", eip_status_text(rc));
        return rc;
    }
    rc = eip_tag_read(&s, "Program:Main.Speed", 1, buf, &len, &type, &status);
    if (rc == EIP_OK && type == EIP_DINT) {
        printf("Speed = %d\n", (int)(int32_t)eip_load_le32(buf));
    } else if (rc == EIP_ERR_CIP) {
        printf("CIP status 0x%02X / 0x%04X\n",
               eip_last_cip_status(&s), eip_last_cip_extended(&s));
    }
    eip_session_close(&s);
    return rc;
}
\end{lstlisting}

\subsection{Rules for Sessions}

\begin{itemize}
  \item \code{eip_session_open} writes over the session. Close a session
        before you open it again.
  \item \code{eip_session_is_open} is safe on a zeroed session that you
        never opened.
  \item After \code{EIP_ERR_TIMEOUT}, the session usually stays open. The
        next request discards a late reply.
  \item A packet that is too short or too large closes the session, because
        the stream cannot be synchronized again. The call that failed returns
        its error. Later calls return \code{EIP_ERR_STATE} until you open the
        session again.
  \item \code{eip_session_is_open} tells you only that the session is
        registered here. It stays 1 after the peer stops.
\end{itemize}

\begin{caution}
Do not use \code{eip_session_is_open} to find a lost device. After
\code{EIP_ERR_IO}, or when requests continue to time out, close the session
and open it again.
\end{caution}

\section{Tags}\label{sec:tags}

\code{eip_tag_read} and \code{eip_tag_write} access Logix tags by name.

\begin{itemize}
  \item Use a period for a member and \code{[n]} or \code{[n,m]} for an
        element: \code{Program:Main.Speed}, \code{Data[2,3].Value}.
  \item The scanner breaks reads and writes that are larger than one packet
        into fragments (Read Tag Fragmented and Write Tag Fragmented).
  \item Data is raw little-endian bytes. Use the helpers
        \code{eip_load_le16}, \code{eip_load_le32}, \code{eip_store_le16}, and
        \code{eip_store_le32}.
  \item \code{cip_type} receives the CIP type code of the tag, for example
        \code{EIP_DINT} (\hex{00C4}). Table~\ref{tab:ciptypes} lists the codes.
\end{itemize}

\begin{lstlisting}
uint8_t buf[4];
uint8_t status = 0;

eip_store_le32(buf, 1234);
eip_tag_write(&s, "Setpoint", EIP_DINT, 1, buf, 4, &status);
\end{lstlisting}

\begin{caution}
A write of more than one fragment is not atomic. If an error occurs during
the write, the PLC can keep the fragments that it already received. Thus the
tag can contain a mix of old and new data. Do not write large arrays that a
running PLC program uses for control while the write is in progress.
\end{caution}

\section{Explicit Messages}\label{sec:explicit}

\code{eip_cip} sends one CIP service with a raw path. Use it for a service that
has no helper. The helpers \code{eip_get_attribute} and
\code{eip_set_attribute} send Get\_Attribute\_Single and
Set\_Attribute\_Single.

\begin{lstlisting}
uint8_t buf[4];
uint16_t len = sizeof buf;
uint8_t status = 0;

/* Identity object (class 0x01), instance 1, attribute 1: vendor ID */
eip_get_attribute(&s, 0x01, 1, 1, buf, &len, &status);
\end{lstlisting}

\begin{itemize}
  \item The encoded request has a limit of 480 bytes.
  \item A general status other than 0 gives \code{EIP_ERR_CIP}. The
        exception is \hex{06} (partial transfer), which gives \code{EIP_OK}
        with the status in \code{*cip_status}.
  \item When a call returns \code{EIP_ERR_CIP}, \code{eip_last_cip_status}
        gives the general status and \code{eip_last_cip_extended} gives the
        extended status.
  \item On a routed session (slot 0 or more), the router can cause the error.
        For example, \hex{01} / \hex{0204} means that the slot did not answer.
\end{itemize}

\section{Discovery and Symbols}\label{sec:discovery}

\begin{lstlisting}
EipDevice found[16];
uint16_t n = 16;
eip_scan_devices(found, &n, 0);       /* broadcast on each interface */

EipIdentity id;
eip_list_identity("192.168.1.10", &id, 0);

static EipSymbol tags[256];
uint16_t count = 256;
eip_list_symbols(&s, tags, &count);   /* Micro800: open with slot -1 */
\end{lstlisting}

\begin{itemize}
  \item \code{eip_scan_devices} sends a broadcast ListIdentity on each
        interface and on the global broadcast address. With a timeout of 0, it
        waits approximately 3~s, sends again each second, and stops 0.5~s
        after the last reply. \code{EIP_OK} with a count of 0 means that no
        device answered.
  \item Broadcasts do not cross routers or VLANs. For a device on a different
        subnet, use \code{eip_list_identity} with its address.
  \item \code{eip_list_symbols} reads the Symbol object. If a later page
        fails, the count still tells you how many symbols the call stored.
  \item \code{eip_device_kind} sorts Rockwell products into Micro800,
        CompactLogix, or other by the catalog name. This is a heuristic.
        Your application decides what to do with the result.
\end{itemize}

\section{Class 1 I/O as an Originator}\label{sec:originator}

The scanner can open a Class~1 connection to an adapter. Your program then
produces \OT{} data and consumes \TO{} data.

\begin{procedure}{Open a Class 1 connection}
\begin{enumerate}
  \item Open a session with \code{slot = -1} and your vendor ID in
        \code{vendor_id}.
  \item Set an \code{EipIoConfig} to zero.
  \item Set \code{config_instance} to the value in the EDS file of the target.
  \item Set \code{output_instance} and \code{output_bytes} for the data that
        you produce.
  \item Set \code{input_instance} and \code{input_bytes} for the data that
        you consume.
  \item Set \code{o2t_rpi_us} and \code{t2o_rpi_us}.
  \item Set \code{run_idle} to 1 if the target expects the run/idle header.
  \item Set \code{timeout_multiplier}.
  \item Call \code{eip_io_open}.
  \item Make sure that the result is \code{EIP_OK}.
\end{enumerate}
\end{procedure}

\begin{lstlisting}
static EipIo io;
EipIoConfig ioc = {0};
uint8_t out[32] = {0}, in[32];

ioc.config_instance = 151;
ioc.output_instance = 150;  ioc.output_bytes = 32;   /* we produce (O->T) */
ioc.input_instance  = 100;  ioc.input_bytes  = 32;   /* we consume (T->O) */
ioc.o2t_rpi_us = ioc.t2o_rpi_us = 10000;             /* 10 ms */
ioc.run_idle = 1;
ioc.timeout_multiplier = 2;                          /* timeout = RPI x 16 */

if (eip_io_open(&s, &io, &ioc) == EIP_OK) {
    for (;;) {                                       /* once per RPI */
        uint16_t in_len = sizeof in;
        eip_io_produce(&io, out, 32, 1 /* run */);
        eip_io_consume(&io, in, &in_len, 100);
    }
    eip_io_close(&s, &io);
}
\end{lstlisting}

These rules apply:

\begin{itemize}
  \item The originator binds UDP 2223--2231, unless you set
        \code{local_port}. Thus it can share a device with an adapter on 2222.
  \item The library does not set the packet interval. Call
        \code{eip_io_produce} one time each RPI.
  \item \code{eip_io_consume} waits for a packet. A loop like the example
        thus runs at the \TO{} rate. A timeout of 0 means 1~s. There is no
        non-blocking form. Use a short timeout, or use different tasks for
        produce and consume.
  \item \code{config_instance = 0} sends instance 0 in the path. The QuickEIP
        adapter accepts it, but a different target can refuse it.
  \item The connection size sent to the target follows
        equation~\eqref{eq:o2t}.
\end{itemize}

\begin{caution}
Use the address that the adapter reports, which is the address that a scan
shows. The adapter tells the originator to send to that address, and it
accepts \OT{} data only from the address of the Forward\_Open. If you use a
different address of the same host, for example \code{127.0.0.1}, the adapter
drops the packets and the connection times out. Explicit messages are not
affected.
\end{caution}

\begin{note}
\notelabel Open Class~1 on a session with \code{slot = -1}. On a routed
session, the Forward\_Open goes through Unconnected Send. That combination has
not been tested.
\end{note}

\subsection{Recover a Lost Connection}\label{sec:recover}

\code{eip_io_produce} uses UDP, so it continues to succeed after the target
stops. A lost connection shows only in \code{eip_io_consume}, which returns
\code{EIP_ERR_TIMEOUT} when no \TO{} packet arrives. Figure~\ref{fig:recover}
shows the recovery sequence.

\begin{figure}[htbp]
\centering
\begin{tikzpicture}[
  font=\sffamily\small,>=Stealth,node distance=0.55cm,
  stp/.style={draw=steel,fill=white,rounded corners=2pt,minimum width=4.6cm,minimum height=0.72cm,align=center},
  call/.style={stp,fill=codebg,font=\ttfamily\small},
  dec/.style={diamond,aspect=2.6,draw=amber,fill=amber!8,inner sep=1pt,align=center,font=\sffamily\scriptsize},
  lab/.style={font=\sffamily\scriptsize,text=steel}]
  \node[stp,fill=teal!12,draw=teal] (loop) {Produce and consume each RPI};
  \node[dec,below=of loop] (to) {Consume returns\\\code{EIP_ERR_TIMEOUT}?};
  \node[dec,below=of to] (n) {Timeouts in a row\\$\geq N$?};
  \node[call,below=0.7cm of n] (c1) {eip\_io\_close};
  \node[call,below=of c1] (c2) {eip\_session\_close};
  \node[call,below=of c2] (c3) {eip\_session\_open};
  \node[call,below=of c3] (c4) {eip\_io\_open};
  \node[stp,right=1.4cm of c3,minimum width=3.0cm] (wait) {Wait, then\\try again};
  \draw[->,steel] (loop) -- (to);
  \draw[->,steel] (to) -- node[lab,right]{yes} (n);
  \draw[->,steel] (to.east) -- ++(1.6,0) |- node[lab,pos=0.25,right]{no: count = 0} (loop.east);
  \draw[->,steel] (n.west) -- ++(-1.2,0) |- node[lab,pos=0.25,left]{no} (loop.west);
  \draw[->,steel] (n) -- node[lab,right]{yes} (c1);
  \draw[->,steel] (c1) -- (c2);
  \draw[->,steel] (c2) -- (c3);
  \draw[->,steel] (c3) -- (c4);
  \draw[->,steel] (c3.east) -- node[lab,above]{fails} (wait.west);
  \draw[->,steel] (wait.north) |- (c2.east);
  \draw[->,steel] (c4.west) -- ++(-2.2,0) node[lab,above,pos=0.5]{OK} |- (loop.west);
  \draw[->,steel] (c4.east) -| node[lab,below,pos=0.25]{fails} (wait.south);
\end{tikzpicture}
\caption{Recovery of a lost Class~1 connection on the originator. The four
calls are all necessary: a reopen of I/O on the old session fails.}
\label{fig:recover}
\end{figure}

\begin{procedure}{Recover a lost Class 1 connection}
\begin{enumerate}
  \item Count the calls to \code{eip_io_consume} in a row that return
        \code{EIP_ERR_TIMEOUT}.
  \item When the count gets to your limit, call \code{eip_io_close}.
  \item Call \code{eip_session_close}.
  \item Call \code{eip_session_open}.
  \item If the open fails, wait and do step 3 again.
  \item Call \code{eip_io_open}.
  \item Set the count to 0, and start the produce and consume loop again.
\end{enumerate}
\end{procedure}

\begin{note}
\notelabel When the peer is gone, \code{eip_io_close} can block for the session
timeout and return an error. This is the expected result. The call still
releases the UDP socket.
\end{note}

If you do not close and open the session, \code{eip_io_open} returns
\code{EIP_ERR_IO}.
